\section{Literature and technical background}
Natural-gas SMR combines high-temperature reforming with water--gas shift (WGS), hydrogen purification and substantial process heat. Literature and authoritative sources establish three facts central to this study: lifecycle emissions cannot be inferred from stack capture alone because upstream natural-gas emissions remain material; high-temperature gas-cooled reactors can provide high-grade helium heat in principle, but reactor-to-chemical-plant coupling remains an integration problem rather than a demonstrated Singapore deployment; and CCS energy and cost depend strongly on the pressure and composition of the captured stream.

The structured literature matrix in \texttt{literature/RESEARCH\_MATRIX.md} records the evidence basis and contradictions used in model construction. International Energy Agency Greenhouse Gas R\&D Programme (IEAGHG) SMR+CCS cases~\cite{ieaghg2017smr} provide the principal process benchmark; National Institute of Standards and Technology (NIST)/JANAF thermochemistry~\cite{nistwebbook} underpins equilibrium and enthalpy calculations; JAEA / Japan Atomic Energy Research Institute (JAERI) HTTR and GTHTR300C work~\cite{jaeahttr,jaeri2004gthtr300c} provides temperature, IHX and helium-loop precedent; Singapore Government sources~\cite{mti2025nuclear,mti2026inir,mti2025carbon,mti2025indonesiaccs} define nuclear and cross-border CCS as future or conditional infrastructure rather than existing assets.



\subsection{Three-figure orientation: baseline, nuclear power, proposed integration}
Before the detailed derivations, Figures~\ref{fig:smr-explainer}--\ref{fig:proposed-orientation} answer three different questions: how conventional hydrogen production obtains reformer heat; what an HTGR normally does with reactor heat in a power-production pathway; and what this project changes by diverting high-temperature heat to the hydrogen process.

\begin{figure}[htbp]\centering
\resizebox{\textwidth}{!}{\input{../results/final_design/generated/figure1_conventional_smr_ccs.tex}}
\caption{Figure 1---conventional hydrogen-production orientation. SMR = Steam Methane Reformer/Reforming; WGS = Water-Gas Shift; CCS = Carbon Capture and Storage; PSA = Pressure Swing Adsorption; T\&S = carbon-dioxide Transport and Storage. The fired furnace supplies the high-temperature reformer heat service; process gas then passes through WGS, CO$_2$ capture and PSA to hydrogen product, while captured CO$_2$ proceeds to conditioning and T\&S. This is an original static engineering schematic reconstructed from the source process; equipment silhouettes are schematic art rather than design drawings~\cite{ieaghg2017smr,inltev953}.}
\label{fig:smr-explainer}
\end{figure}

\begin{figure}[htbp]\centering
\resizebox{0.94\textwidth}{!}{\input{../results/final_design/generated/figure2_htgr_power.tex}}
\caption{Figure 2---conceptual conventional HTGR electricity-generation pathway. HTGR = High-Temperature Gas-Cooled Reactor; He = helium. Primary helium forms a closed reactor-to-steam-generator loop; the separate steam/water cycle drives a turbine and generator and returns condensate through a pump. This is a reference architecture only and does \textbf{not} assign or imply a numerical electricity output for the proposed CN4252 system. Equipment silhouettes are schematic art, not mechanical design drawings.}
\label{fig:htgr-power-orientation}
\end{figure}

\begin{figure}[p]\centering
\resizebox{\textwidth}{!}{\input{../results/final_design/generated/figure3_proposed_htgr_smr_ccs.tex}}
\caption{Figure 3---proposed HTGR-assisted SMR+CCS system. HTGR = High-Temperature Gas-Cooled Reactor; IHX = Intermediate Heat Exchanger; He = helium; SMR = Steam Methane Reformer; WGS = Water-Gas Shift; PSA = Pressure Swing Adsorption; CCS = Carbon Capture and Storage; T\&S = Transport and Storage. Idaho National Laboratory process evidence anchors the 925$^\circ$C reactor-outlet case, 900$^\circ$C secondary-He supply, approximately 466$^\circ$C return, 78.49 kg/s source flow, 871$^\circ$C reformer outlet, 176.8 MWth duty, S/C=3.0 and 88\% PSA hydrogen recovery~\cite{inltev953,inltev961}. The 600 MWth GTHTR300C-class reactor basis is a literature design basis~\cite{jaeri2004gthtr300c,nishihara2007potential}. Primary and secondary helium are physically separated by the IHX and neither mixes with the SMR process gas. The approximately 170 MWth value is a physical-IHX reference/qualification caveat and does not establish a qualified 176.8 MWth project IHX. The approximately 370/371 MWth value is the source architecture/economic process-heat branch, \textbf{not} one physical IHX~\cite{nishihara2007potential}. The 423.2 MWth value is remaining reactor thermal capacity, \textbf{not} electricity output. The 88\% value is PSA hydrogen recovery, \textbf{not} hydrogen purity. H$_2$ product is 130 MMSCFD; annualised output is approximately 97,946 t/y. The CO$_2$ endpoint is a conditional transport/storage screening boundary rather than guaranteed infrastructure.}
\label{fig:proposed-orientation}
\end{figure}

\subsection{How the proposed system works, step by step}
\begin{enumerate}
\item \textbf{Fission produces heat.} Nuclear fission releases energy primarily as fission-fragment kinetic energy and associated energy deposition, heating the reactor fuel and core; subsequent radioactive decay contributes decay heat.
\item \textbf{Primary helium removes core heat.} At the selected process operating point, the source reactor-outlet case is 925$^\circ$C~\cite{inltev961}.
\item \textbf{Primary helium reaches the IHX.} The hot primary coolant carries reactor heat to the Intermediate Heat Exchanger.
\item \textbf{The IHX provides physical separation.} Heat crosses exchanger walls; primary reactor helium does not enter the chemical plant.
\item \textbf{Secondary helium receives the heat.} INL Case 6 supplies secondary helium at 900$^\circ$C with a reported process-side flow of 78.49 kg/s~\cite{inltev961}.
\item \textbf{Secondary helium heats the reformer.} The source process requires 176.8 MWth and reports an 871$^\circ$C reformer outlet~\cite{inltev953,inltev961}.
\item \textbf{Secondary helium returns.} After heat transfer, the source return state is approximately 466$^\circ$C before reheating at the IHX~\cite{inltev961}.
\item \textbf{Natural gas and steam enter the reformer.} The selected steam-to-carbon molar ratio is 3.0~\cite{inltev953}. Nuclear heat replaces fossil \emph{process heat}; natural gas remains the chemical hydrogen feedstock.
\item \textbf{Steam methane reforming produces synthesis gas.} Methane reacts with steam to form a hydrogen-rich mixture containing carbon monoxide.
\item \textbf{Water-gas shift increases hydrogen.} Carbon monoxide reacts with steam toward CO$_2$ and additional H$_2$.
\item \textbf{CO$_2$ is separated.} The source process represents amine-based capture of the concentrated carbon stream.
\item \textbf{PSA recovers hydrogen.} The selected Pressure Swing Adsorption recovery is 88\%~\cite{inltev953}.
\item \textbf{Hydrogen leaves as product.} The source service is 130 million standard cubic feet per day (MMSCFD) H$_2$~\cite{inltev953,inltev961}; annualised at 85\% availability, the project calculation gives approximately 97,946 t H$_2$/y.
\item \textbf{Captured CO$_2$ is conditioned.} The source captured stream is 1,927 short ton/day~\cite{inltev961}, annualised to approximately 542,362 t/y at the declared availability.
\item \textbf{CO$_2$ leaves the project boundary.} Compression/conditioning leads to transport and geological storage; the cross-border route and storage site remain conditional rather than demonstrated.
\end{enumerate}

\subsection{Visual orientation: from Singapore context to the process}
Figure~\ref{fig:singapore-context} gives the external-data context before introducing the process model. Singapore's 2024 electricity system remained strongly gas-dependent: natural gas supplied 94.0\% of the generation fuel mix and the average grid emission factor was 0.402 kgCO$_2$/kWh~\cite{ema2024context}; Singapore's Energy Market Authority (EMA) also describes power generation as contributing about 40\% of national carbon emissions~\cite{ema2024futurefund}. These are external Singapore statistics, not outputs of this project.

\begin{figure}[htbp]
\centering
\begin{tikzpicture}[box/.style={draw,rounded corners,align=center,minimum width=0.27\linewidth,minimum height=22mm}]
\node[box] (a) {\Large\textbf{94.0\%}\\[1mm]Natural gas share\\of 2024 electricity fuel mix};
\node[box,right=8mm of a] (b) {\Large\textbf{0.402}\\[1mm]kgCO$_2$/kWh\\2024 grid factor};
\node[box,right=8mm of b] (c) {\Large\textbf{$\sim$40\%}\\[1mm]of national carbon emissions\\from power generation};
\end{tikzpicture}
\caption{Singapore energy context from Energy Market Authority (EMA) 2024 data~\cite{ema2024context,ema2024futurefund}. The panel is reconstructed from authoritative reported values stored in \texttt{data/external/ema\_singapore\_energy\_context\_2024.csv}; it is not a project-model result.}
\label{fig:singapore-context}
\end{figure}

Figure~\ref{fig:smr-explainer} introduces the hydrogen pathway at a deliberately simpler level than the detailed model. Feed carbon enters with natural gas. Reforming converts methane and steam to syngas; WGS converts CO and steam toward CO$_2$ and additional H$_2$; CO$_2$ separation removes a concentrated carbon stream; PSA produces high-purity hydrogen and a carbon-bearing tail gas. CCS changes where much of the carbon goes, but it does not remove the feedstock carbon from the chemistry~\cite{ieaghg2017smr}.


The nuclear proposal changes the \emph{heat source}, not the feed-carbon chemistry. Figure~\ref{fig:nuclear-heat-explainer} separates those ideas. A fired reformer supplies high-temperature heat by combustion; the nuclear-assisted concept instead transfers reactor heat through helium and an intermediate heat exchanger. The final design hierarchy is a 950$^\circ$C-class JAEA primary-helium hardware envelope $\rightarrow$ 900$^\circ$C delivered process heat $\rightarrow$ 871$^\circ$C reformer outlet. The specific INL Case-6 ROT is 925$^\circ$C~\cite{inltev961}, while JAEA/JAERI provides the 950$^\circ$C-class hardware precedent~\cite{jaeri2004gthtr300c}. Feedstock carbon still requires capture and lifecycle accounting.

\begin{figure}[htbp]
\centering
\begin{tikzpicture}[>=Latex,node distance=8mm,unit/.style={draw,rounded corners,align=center,minimum height=10mm,minimum width=25mm},arr/.style={->,thick},heat/.style={->,thick,dashed}]
\node[unit] (fire) {Conventional\\fuel combustion};
\node[unit,right=of fire] (r1) {Reformer\\high-temperature heat};
\draw[heat](fire)--node[above,font=\scriptsize]{fired heat}(r1);
\node[below=13mm of fire,unit] (reactor) {HTGR-class reactor\\primary He 950$^\circ$C};
\node[right=of reactor,unit] (ihx) {IHX\\process heat 900$^\circ$C};
\node[right=of ihx,unit] (r2) {Reformer\\gas outlet 871$^\circ$C};
\draw[heat](reactor)--(ihx);\draw[heat](ihx)--(r2);
\node[right=10mm of r1,align=left] {\textbf{What changes:} heat service\\\textbf{What does not:} feed carbon};
\end{tikzpicture}
\caption{Original two-layer explanation of nuclear assistance. The concept substitutes a high-temperature heat-delivery chain for fired reformer heat; it does not eliminate carbon entering as natural-gas feed. Temperatures show the final literature-supported design ladder; the 950$^\circ$C value is the JAEA hardware envelope and the 900/871$^\circ$C process pair is anchored to INL~\cite{inltev961,jaeri2004gthtr300c}.}
\label{fig:nuclear-heat-explainer}
\end{figure}


\subsection{Why an HTGR can supply process heat safely across an interface}
A high-temperature gas-cooled reactor (HTGR) uses helium as coolant. Helium is chemically inert and remains gaseous at reactor temperatures, making it suitable for carrying high-grade heat. HTGR fuel commonly uses coated particles (TRISO): the fuel kernel is surrounded by multiple ceramic/carbon barrier layers. For this project, the important point is not reactor-physics detail but the resulting engineering capability to operate with high-temperature helium~\cite{jaeahttr,jaeri2004gthtr300c}.

The intermediate heat exchanger (IHX) is the isolation boundary between the reactor primary circuit and the chemical heat-transfer circuit. Primary helium transfers heat across exchanger walls to a separate secondary-helium loop; the secondary helium then heats the reformer. Thus \textbf{reactor primary coolant does not enter the hydrogen process}.

\begin{figure}[htbp]\centering
\begin{tikzpicture}[>=Latex,node distance=10mm,unit/.style={draw,rounded corners,align=center,minimum height=13mm,minimum width=29mm},arr/.style={->,thick},heat/.style={->,thick,dashed}]
\node[unit] (core) {HTGR core\\TRISO fuel\\primary helium};
\node[unit,right=of core] (ihx) {IHX\\metal boundary\\heat crosses};
\node[unit,right=of ihx] (sec) {Secondary helium\\separate loop\\900$^\circ$C supply};
\node[unit,right=of sec] (ref) {Chemical reformer\\176.8 MWth\\871$^\circ$C outlet};
\draw[heat](core)--node[above,font=\scriptsize]{primary He}(ihx);
\draw[heat](ihx)--node[above,font=\scriptsize]{secondary He}(sec);
\draw[heat](sec)--(ref);
\node[draw,dashed,fit=(core),inner sep=5mm,label=below:{\scriptsize nuclear primary boundary}] {};
\node[draw,dashed,fit=(sec)(ref),inner sep=5mm,label=below:{\scriptsize non-nuclear process-heat side}] {};
\end{tikzpicture}
\caption{Original isolation schematic reconstructed from HTGR/IHX architecture reported by JAEA/JAERI~\cite{jaeahttr,jaeri2004gthtr300c,nishihara2007potential}. It emphasizes the safety-relevant conceptual point: heat crosses the IHX, while the primary reactor coolant and chemical process remain separate circuits.}
\label{fig:ihx-isolation}
\end{figure}

The user-provided project-source diagrams were inspected as technical visual references during final QA. They contain useful depictions of fired reforming, amine capture, PSA, HTGR helium loops and IHX coupling, but their supplied filenames do not consistently describe their pictured systems. They were therefore not copied into the manuscript; the explanatory figures above were recreated from source-supported engineering concepts with primary-source attribution.
